\subsection{Tensor products and exterior powers of sesquilinear forms.}

In this section, A denotes a commutative ring. A bilinear form on a product of two A-modules is therefore a special case of a sesquilinear form. We denote by J an automorphism of A, and by J' its inverse.

Let \(\mathbf{E}_i\) ( \(i = 1, \ldots, m\) ) be A-modules. The map

\[
(x _ {1}, \dots , x _ {m}) \rightarrow x _ {1} \otimes \dots \otimes x _ {m}
\]

of \(\prod_{i=1}^{m}\mathrm{E}_{i}^{j}\) into \((\bigotimes_{i=1}^{m}\mathrm{E}_{i})^{j}\) ( \(x_{i}\in\mathrm{E}_{i}^{j}\) ) (cf. def. 5, no. 2) is evidently A-multilinear; it therefore defines (chap. III, § 1, no. 7) an A-linear map \(f\) of \(\bigotimes_{i}\mathrm{E}_{i}^{\mathrm{J}}\) into \((\bigotimes_{i}\mathrm{E}_{i})^{\mathrm{J}}\) ; this map sends \(x_{1}\otimes\cdots\otimes x_{m}\) (where the symbols \(\otimes\) denote tensor products in \(\bigotimes_{i}\mathrm{E}_{i}^{\mathrm{J}}\) ) to \(x_{1}\otimes\cdots\otimes x_{m}\) (where the symbols \(\otimes\) denote tensor products in \((\bigotimes_{i}\mathrm{E}_{i})^{\mathrm{J}}\) ). Hence \(f\) is an isomorphism of \(\bigotimes_{i}\mathrm{E}_{i}^{\mathrm{J}}\) onto \((\bigotimes_{i}\mathrm{E}_{i})^{\mathrm{J}}\) . We shall identify these two modules by means of this isomorphism.

Similarly, let E be an A-module. The map

\[
(x _ {1}, \dots , x _ {m}) \rightarrow x _ {1} \wedge \dots \wedge x _ {m}
\]

of \((\mathrm{E}^{\mathfrak{J}})^{m}\) into \((\wedge^{\mathfrak{m}}\mathrm{E})^{\mathfrak{J}}\) is obviously A-multilinear and alternating. It therefore defines an A-linear map \(f\) of \(\wedge^{\mathfrak{m}}\mathrm{E}^{\mathfrak{J}}\) into \((\wedge^{\mathfrak{m}}\mathrm{E})^{\mathfrak{J}}\) , which is evidently an isomorphism. We shall identify \(\wedge^{\mathfrak{m}}\mathrm{E}^{\mathfrak{J}}\) and \((\wedge^{\mathfrak{m}}\mathrm{E})^{\mathfrak{J}}\) by means of this isomorphism.

Let \(x'\) be an element of the dual \(\mathbf{E}^*\) of \(\mathbf{E}\) . The map \(x \to \langle x, x' \rangle^{\mathfrak{J}}\) ( \(x \in \mathbf{E}\) ) is an element \(x'^{\mathfrak{J}}\) of \((\mathrm{E}^{\mathfrak{J}})^*\) , and it is immediate that \(x' \to x'^{\mathfrak{J}}\) is a bijection \(g\) of \(\mathbf{E}^*\) onto \((\mathrm{E}^{\mathfrak{J}})^*\) satisfying \(g(ax') = a^{\mathfrak{J}}g(x')\) for all \(a \in \mathbf{A}\) . Consequently the composite map of \(g\) and of the identity map of \((\mathrm{E}^*)^{\mathfrak{J}}\) onto \(\mathbf{E}^*\) is an isomorphism of \((\mathrm{E}^*)^{\mathfrak{J}}\) onto \((\mathrm{E}^{\mathfrak{J}})^*\) . We shall identify these modules by means of this isomorphism, and we shall denote them \(\mathbf{E}_j^*\) .

Let \(E_{i}, F_{i} (i = 1, \ldots, m)\) be A-modules, and \(\Phi_{i} (i = 1, \ldots, m)\) a sesquilinear form for J on \(E_{i} \times F_{i}\) . The map

\[
(x _ {1}, \dots , x _ {m}, y _ {1}, \dots , y _ {m}) \rightarrow \Phi_ {1} (x _ {1}, y _ {1}) \Phi_ {2} (x _ {2}, y _ {2}) \dots \Phi_ {m} (x _ {m}, y _ {m})
\]

\((x_{i} \in \mathrm{E}_{i}, y_{i} \in \mathrm{F}_{i}, i = 1, \ldots, m)\) is an A-multilinear map from \(\mathrm{E}_{1} \times \cdots \times \mathrm{E}_{m} \times \mathrm{F}_{1}^{J} \times \cdots \times \mathrm{F}_{m}^{J}\) into A, and thus defines a bilinear form on \((\bigotimes_{i} \mathrm{E}_{i}) \times (\bigotimes_{i} \mathrm{F}_{i}^{J})\) (chap. III, § 1, no. 7). Since the second factor has been identified with \((\bigotimes_{i} \mathrm{F}_{i})^{J}\) , we have therefore defined a sesquilinear form \(\Phi\) for J on \((\bigotimes_{i} \mathrm{E}_{i}) \times (\bigotimes_{i} \mathrm{F}_{i})\) . It is characterized by

\[
\Phi (x _ {1} \otimes \dots \otimes x _ {m}, y _ {1} \otimes \dots \otimes y _ {m}) = \prod_ {i = 1} ^ {m} \Phi_ {i} (x _ {i}, y _ {i}) \quad (x _ {i} \in \mathrm{E} _ {i}, y _ {i} \in \mathrm{F} _ {i}).\tag{35}
\]

DEFINITION 11. --- Given A-modules \(\mathbf{E}_{i}, \mathbf{F}_{i}\) ( \(i=1,\ldots,m\) ) and, for each \(i\) , a sesquilinear form \(\Phi_{i}\) for J on \(\mathbf{E}_{i} \times \mathbf{F}_{i}\) , the sesquilinear form \(\Phi\) for J on \((\bigotimes_{i}\mathbf{E}_{i}) \times (\bigotimes_{i}\mathbf{F}_{i})\) characterized by (35) is called the tensor product of the sesquilinear forms \(\Phi_{i}\) .

In the case where the \(E_{i}\) and the \(F_{i}\) are equal to a single module E, and where the \(\Phi_{i}\) are equal to a single form \(\Psi\) , one says that \(\Phi\) is the extension of \(\Psi\) to \(\otimes^{m}E\) .

With the notations of def. 11, let us study the maps associated with \(\Phi\) . From formula (24) (no. 6) and (35) we derive the relation

\[
\Phi (x _ {1} \otimes \dots \otimes x _ {m}, y _ {1} \otimes \dots \otimes y _ {m}) = \prod_ {i = 1} ^ {m} \left\langle x _ {i}, d _ {\Phi_ {i}} (y _ {i}) \right\rangle = \prod_ {i = 1} ^ {m} \left\langle y _ {i}, s _ {\Phi_ {i}} (x _ {i}) \right\rangle^ {j}.
\]

We therefore have:

\[
s _ {\Phi} = j _ {s} \circ (s _ {\Phi_ {1}} \otimes \dots \otimes s _ {\Phi_ {m}}), \quad d _ {\Phi} = j _ {d} \circ (d _ {\Phi_ {1}} \otimes \dots \otimes d _ {\Phi_ {m}})\tag{36}
\]

where \(j_s\) (resp. \(j_d\) ) denotes the canonical map of \(\bigotimes_i F_i^*\) into \((\bigotimes_i F_i)^*\) (resp. of \(\bigotimes_i E_i^*\) into \((\bigotimes_i E_i)^*\) ) (chap. III, § 1, nos. 4 and 7).

PROPOSITION 9. — Let A be a commutative field equipped with an automorphism J, E \(_{i}\) , F \(_{i}\) two finite-dimensional vector spaces over A, and Φ \(_{i}\) a sesquilinear form for J on E \(_{i}\) × F \(_{i}\) (1 ≤ i ≤ m). If the forms Φ \(_{i}\) are nondegenerate, then so is their tensor product Φ. In this case the inverse form \(\hat{\Phi}\) of \(\Phi\) is the tensor product of the inverse forms Φ \(_{i}\) .

Indeed, since A is a field, it follows from Props. 6 and 7 of Chap. III, § 1, no. 3 that a tensor product of injective (resp. surjective) linear maps of A-modules is an injective (resp. surjective) linear map. Since the \(s_{\Phi_i}\) are bijective by hypothesis (prop. 6, no. 6), hence so is their tensor product. On the other hand, the canonical map \(j_s\) of \(\bigotimes_i F_i^*\) into \((\bigotimes_i F_i)^*\) is bijective (chap. III, § 1, no. 5, prop. 11). Therefore, by virtue of (36), \(s_\Phi\) is bijective, and this establishes our first assertion (prop. 6, no. 6). Likewise \(d_\Phi\) is bijective.

In the second assertion we have implicitly identified \(\bigotimes_{i} \mathrm{F}_{i}^{*}\) with \((\bigotimes_{i} \mathrm{F}_{i})^{*}\) and \(\bigotimes_{i} \mathrm{E}_{i}^{*}\) with \((\bigotimes_{i} \mathrm{E}_{i})^{*}\) by means of the maps \(j_{s}\) and \(j_{d}\) , which are here isomorphisms. The inverse forms cited in the statement exist since the \(s_{\Phi_{i}}\) , the \(d_{\Phi_{i}}\) , \(s_{\Phi}\) and \(d_{\Phi}\) are bijective (no. 7). Then set \(x' = x_{1}' \otimes \cdots \otimes x_{m}'\) , \(y' = y_{1}' \otimes \cdots \otimes y_{m}'\) ( \(x_{i}' \in \mathrm{E}_{i}^{*}\) , \(y_{i}' \in \mathrm{F}_{i}^{*}\) , \(i = 1, \ldots, m\) ). By definition of the inverse forms, and in view of (36), we have

\[
\begin{array}{l} \widehat {\Phi} (j _ {s} (y ^ {\prime}), j _ {d} (x ^ {\prime})) = \Phi (s _ {\Phi_ {1}} ^ {- 1} (y _ {1} ^ {\prime}) \otimes \dots \otimes s _ {\Phi_ {m}} ^ {- 1} (y _ {m} ^ {\prime}), d _ {\Phi_ {1}} ^ {- 1} (x _ {1} ^ {\prime}) \otimes \dots \otimes d _ {\Phi_ {m}} ^ {- 1} (x _ {m} ^ {\prime})) \\ = \prod_ {i = 1} ^ {m} \Phi_ {i} (s _ {\Phi_ {i}} ^ {- 1} (y _ {i} ^ {\prime}), d _ {\Phi_ {i}} ^ {- 1} (x _ {i} ^ {\prime})) = \prod_ {i = 1} ^ {m} \widehat {\Phi} _ {i} (y _ {i} ^ {\prime}, x _ {i} ^ {\prime}), \end{array}
\]

whence our second assertion.

QED.

Let E and F be two modules over the commutative ring A, and \(\Phi\) a sesquilinear form for J on E × F. The map

\[
(x _ {1}, \dots , x _ {m}, y _ {1}, \dots , y _ {m}) \rightarrow \det (\Phi (x _ {i}, y _ {k})) \quad (x _ {i} \in \mathrm{E}, y _ {i} \in \mathrm{F}, i = 1, \dots , m)
\]

of \(\mathbf{E}^{m}\times(\mathbf{F}^{\mathrm{j}})^{m}\) into A is A-multilinear. It therefore defines a bilinear form \(\Phi'\) on \((\bigotimes^{\overline{m}}\mathbf{E})\times(\bigotimes^{\overline{m}}\mathbf{F}^{\mathrm{j}})\) characterized by

\[
\Phi^ {\prime} (x _ {1} \otimes \dots \otimes x _ {m}, y _ {1} \otimes \dots \otimes y _ {m}) = \det (\Phi (x _ {i}, y _ {k})).
\]

Since the first member is zero when \(x_{i}=x_{k}\) or \(y_{i}=y_{k}\) ( \(i\neq k\) ), \(\Phi^{\prime}\) defines, by passing to quotients, a bilinear form on \((\bigwedge^{m}\mathrm{E})\times(\bigwedge^{m}\mathrm{F}^{\mathrm{J}})\) or again, since \(\bigwedge^{m}\mathrm{F}^{\mathrm{J}}\) is identified with \((\bigwedge^{m}\mathrm{F})^{\mathrm{J}}\) , a form \(\Phi_{(m)}\) sesquilinear for J on \((\bigwedge^{m}\mathrm{E})\times(\bigwedge^{m}\mathrm{F})\) . It is characterized by

\[
\left\{ \begin{array}{l} \Phi_ {(m)} (x _ {1} \wedge \dots \wedge x _ {m}, y _ {1} \wedge \dots \wedge y _ {m}) = \det (\Phi (x _ {i}, y _ {k})) \\ (x _ {i} \in \mathrm{E}, y _ {i} \in \mathrm{F}, i = 1, \dots , m). \end{array} \right.\tag{37}
\]

DEFINITION 12. --- Given two A-modules E, F and a form \(\Phi\) sesquilinear for J on \(E \times F\) , the form \(\Phi_{(m)}\) sesquilinear for J on \((\bigwedge^{m} E) \times (\bigwedge^{m} F)\) characterized by (37) is called the extension of \(\Phi\) to the m-th exterior powers.

With the notation of def. 12, let us study the maps associated with \(\Phi_{(m)}\) . From formula (24) (no. 6) and (37) we derive the relations

\[
\begin{array}{r l} \Phi_ {(m)} (x _ {1} \wedge \dots \wedge x _ {m}, y _ {1} \wedge \dots \wedge y _ {m}) & = \det (\langle x _ {i}, d _ {\Phi} (y _ {k}) \rangle) \\ & = \det (\langle y _ {i}, s _ {\Phi} (x _ {k}) \rangle^ {j}). \end{array}
\]

We therefore have

\[
s _ {\Phi (m)} = k _ {s} \circ (\stackrel {m} {\wedge} s _ {\Phi}), \quad d _ {\Phi_ {m}} = k _ {d} \circ (\stackrel {m} {\wedge} d _ {\Phi}),\tag{38}
\]

where \(k_s\) (resp. \(k_d\) ) denotes the canonical map of \(\bigwedge^m F^*\) into \((\bigwedge^m F)^*\) (resp. of \(\bigwedge^m E^*\) into \((\bigwedge^m E)^*\) ) (cf. Chap. III, § 8, no. 2).

PROPOSITION 10. --- Let A be a commutative field equipped with an automorphism J, E and F two finite-dimensional vector spaces over A, and \(\Phi\) a sesquilinear form for J on E \(\times\) F. If \(\Phi\) is nondegenerate, then its extension \(\Phi_{(m)}\) to the m-th exterior powers is nondegenerate, and the inverse form of \(\Phi_{(m)}\) is the extension to the m-th exterior powers of the inverse form \(\hat{\Phi}\) of \(\Phi\) .

Indeed, as \(s_{\Phi}\) and \(d_{\Phi}\) are bijective by hypothesis (prop. 6, no. 6), and the same holds for their exterior powers (chap. III, § 5, no. 7). On the other hand, the canonical maps \(k_s\) and \(k_d\) are bijective (chap. III, § 8, no. 2, th. 1). Therefore, by virtue of (38), \(s_{\Phi(m)}\) and \(d_{\Phi(m)}\) are bijective, which proves that \(\Phi_{(m)}\) is nondegenerate (prop. 6, no. 6). In the second assertion we have implicitly identified \(\bigwedge^m F^*\) with \((\bigwedge^m F)^*\) and \(\bigwedge^m E^*\) with \((\bigwedge^m E)^*\) by means of the maps \(k_s\) and \(k_d\) , which are here isomorphisms (loc. cit.). The inverse forms considered in the statement exist since \(s_{\Phi}\) , \(d_{\Phi}\) , \(s_{\Phi(m)}\) and \(d_{\Phi(m)}\) are bijective (no. 7). Then set \(x' = x_1' \wedge \ldots \wedge x_m'\) and \(y' = y_1' \wedge \ldots \wedge y_m'\) ( \(x_i' \in E^*\) , \(y_i' \in F^*\) , \(i = 1, \ldots, m\) ). By definition of inverse forms (no. 7) and in view of (38), we have

\[
\begin{array}{r l} \hat {\Phi} _ {(m)} (k _ {s} (y ^ {\prime}), k _ {d} (x ^ {\prime})) & = \Phi_ {(m)} (s _ {\Phi} ^ {- 1} (y _ {1} ^ {\prime}) \wedge \dots \wedge s _ {\Phi} ^ {- 1} (y _ {m} ^ {\prime}), d _ {\Phi} ^ {- 1} (x _ {1} ^ {\prime}) \wedge \dots \wedge d _ {\Phi} ^ {- 1} (x _ {m} ^ {\prime})) \\ & = \det (\Phi (s _ {\Phi} ^ {- 1} (y _ {i} ^ {\prime}), d _ {\Phi} ^ {- 1} (x _ {k} ^ {\prime}))) = \det (\hat {\Phi} (y _ {i} ^ {\prime}, x _ {k} ^ {\prime})) \end{array}
\]

whence our second assertion.

Remark. --- Let E be a free A-module, and let θ be the canonical isomorphism from ∧E onto the submodule of antisymmetrized tensors of order m (chap. III, § 5, no. 6, prop. 6). Let Φ be a sesquilinear form on E, let Φ(m) be the extension of Φ to ∧E, and let Θ be the sesquilinear form on ∧E which is the inverse image with respect to θ of the extension of Φ to ⊗E. By the definition of θ and of the antisymmetrization of a tensor, and by (35), we have

\[
\Theta (x _ {1} \wedge \dots \wedge x _ {m}, y _ {1} \wedge \dots \wedge y _ {m}) = \sum_ {\sigma , \tau} \varepsilon_ {\sigma} \varepsilon_ {\tau} \Phi (x _ {\sigma (1)}, y _ {\tau (1)}) \dots \Phi (x _ {\sigma (m)}, y _ {\tau (m)})
\]

where \(\sigma\) and \(\tau\) range over the symmetric group \(\mathfrak{S}_{m}\) . By the formula for computing determinants and formula (37), this expression can be written

\[
\sum_ {\tau \in \mathfrak {S} _ {m}} \varepsilon_ {\tau} \det (\Phi (x _ {i}, y _ {\tau (k)})) = m! \det (\Phi (x _ {i}, y _ {k}));
\]

In other words, we have \(\Theta = m!\Phi_{(m)}\) .

